\section{Two categorical constructions with finite witnesses}
\label{sec:categorycases}

The repository's categorical packet constructs finite limits and glues
ordinary function tables on a finite cover. Its object-calculus development
also derives coherent points for towers with bounded surjective lifting.
We construct the sheaf associated to a supplied presheaf and extract the part of an inverse tower on which lifting works.
Background on finite-space sheaves and stabilization appears in \href{https://arxiv.org/abs/math/0610685}{Ladkani, \S2.2}, and the
\href{https://stacks.math.columbia.edu/tag/0594}{Stacks Project, Tag 0594}. See also the \repo{docs/proofs/foundations/ipra_categorical_methods/paper/sections/05_descent.tex}{ProveIt descent packet}.

\subsection{Sheafification on an arbitrary finite topological space}

Fix a topology on $X=[d]=\{0,\ldots,d-1\}$. Encode an open by its
$d$-bit mask, and the topology by its table on $2^d$ masks.
Require the empty and full masks to be open and check binary unions and
intersections. These suffice because there are finitely many masks.
A finite presheaf $F$ is a number $f_U$ for each open, with
$F(U)=[f_U]$, and a restriction table $r^F_{UV}:F(U)\to F(V)$
for each open $V\subseteq U$. Test identity and composition entrywise.
A morphism of presheaves is a family of tables commuting with restrictions.
A sheaf $G$ is required to pass a finite gluing test. For every cover of $U$
and every family of sections agreeing on pairwise intersections, there is
exactly one section in $G(U)$ restricting to that family. A cover is
a subfamily of the topology whose union is $U$. Removing repeated masks
causes no loss, since matching sections on repeated patches are equal.
This includes the empty cover of the empty open.

For $x\in X$, let $N_x$ be the intersection of all opens containing $x$.
It is open, contains $x$, and is contained in any open containing $x$.
In particular $y\in N_x$ implies $N_y\subseteq N_x$.
For an open $U$, define $SF(U)$ to be the compatible tuples
\[
 (a_x)_{x\in U},\qquad a_x\in F(N_x),\qquad
 r^F_{N_xN_y}(a_x)=a_y\quad(x\in U,\ y\in N_x).
\]
Order all tuples lexicographically and code a compatible tuple by its rank
in the retained list. For the empty open there is one tuple.
Restrict by deleting coordinates. On a morphism $h:F\to G$, apply
$h_{N_x}$ coordinatewise. Naturality of $h$ proves compatibility.
Identity and composition follow by evaluating at each coordinate. Define
$\eta_{F,U}(t)_x=r^F_{UN_x}(t)$. The presheaf identities prove that these
tuples are compatible, natural in $U$, and natural in $F$.

\paragraph{Theorem.}
$S$ is a uniform PR sheafification for finite-set-valued presheaves on
finite spaces. Composition with $\eta_F$ gives a PR bijection
\[
 \operatorname{Hom}_{\rm Sh}(SF,G)\ \longrightarrow \operatorname{Hom}_{\rm PSh}(F,G)
\]
for every supplied finite sheaf $G$. Both inverse equations have
finite arithmetic proofs.

\paragraph{Proof.}
Take matching tuples on a cover of $U$. For each $x\in U$ read its
entry in the first patch containing $x$. Agreement on intersections
makes this independent of the patch. If $y\in N_x$, the chosen patch
containing $x$ also contains $y$, so compatibility of its tuple gives
the required restriction equation. These forced entries prove
existence and uniqueness of the glued tuple. It includes the empty cover.

Over $N_x$, evaluation at $x$ identifies $SF(N_x)$ with $F(N_x)$.
Indeed $b\in F(N_x)$ gives the tuple $y\mapsto r^F_{N_xN_y}(b)$
and compatibility forces every tuple on $N_x$ to have this form.

Given $h:F\to G$ and $a\in SF(U)$, form $g_x=h_{N_x}(a_x)$.
To glue $g_x$ and $g_z$ on $W=N_x\cap N_z$, take any $y\in W$.
Both restrict to $h_{N_y}(a_y)$ on $N_y$, by naturality and
compatibility. The opens $N_y$ cover $W$. Uniqueness in the sheaf
condition makes the two restrictions equal on $W$, including when
$W$ is empty. Thus the $g_x$ match. Define $\bar h_U(a)$ by scanning
$G(U)$ for their unique glue. Its restrictions to smaller opens
give the same answer by uniqueness on their minimal-neighborhood covers.
This proves naturality in $U$. For $t\in F(U)$ the required glue of the
components of $\eta_F(t)$ is $h_U(t)$, proving $\bar h\eta_F=h$.
Conversely, if $k\eta_F=h$, then on $N_x$ the bijection above forces
$k_{N_x}(a|_{N_x})=h_{N_x}(a_x)$. Naturality and uniqueness on
the cover of $U$ give $k_U(a)=\bar h_U(a)$. This proves the second
inverse equation. Compatibility with pre- and postcomposition follows since
the two candidate extensions have the same composite with the unit.
Taking $F$ itself to be a sheaf, its neighborhood sections glue uniquely,
so $\eta_F$ is an isomorphism. In particular $S(SF)\cong SF$.

\paragraph{Arithmetic bounds and example.}
With $q=2^d$ and $D=\max(1,\max_U f_U)$ (the maximum over opens),
there are at most $D^d$ candidate tuples over each open, including the
single empty tuple. If any required coordinate has an empty alphabet, reject the tuple before digit extraction. Mask scans have bound $q$, subfamily
scans $2^q$. If validating a sheaf with fiber bound $H\ge1$, cover
assignments number at most $H^q$, and the uniqueness scan has bound $H$.
Keep the tuple, prefix agreement bits, number of matches truncated at two,
and first matching index. Induct on their bounded prefix specifications.
All induced names, restrictions, unit and extension tables thus use only
PR loops, also when the input topology varies. For the discrete
two-point space, take the constant presheaf $F(U)=[2]$ with identity
restrictions, including at the empty open. Then $SF(X)$ has four sections,
$SF(\varnothing)$ has one, and $\eta_{F,X}$ sends $b$ to $(b,b)$.
This construction applies to spaces with smallest neighborhoods. An arbitrary infinite space may lack
smallest open neighborhoods.

\subsection{Stable images in an inverse tower}

Fix PR bounds $b_n$, Boolean membership predicates for
$A_n\subseteq[b_n]$, and PR bonding maps
$p_n:A_{n+1}\to A_n$ with proofs of typing.
For $m\ge n$ let $P_{n,m}:A_m\to A_n$ denote the composite, identity
when $m=n$. Starting with $z$ at $m$, iterate $j\mapsto
p_{m-1-j}$ on the register for $m-n$ steps.
Splitting this finite loop at any intermediate $k$ gives
$P_{n,k}P_{k,m}=P_{n,m}$. More explicitly, resume the run after
$m-k$ steps. Induction on the number of resumed steps gives identical
registers in both runs because both use the same next bond.
Let $I_{n,m}=P_{n,m}(A_m)$. Membership is one bounded search below $b_m$.
Images decrease as $m$ increases, by the composite law.

A proposed rate is a fixed PR $s(n)\ge n$. Set $B_n=I_{n,s(n)}$,
$K_n=\max(s(n),s(n+1))$. Consider the single comparison at each $n$
\begin{equation}
       I_{n,s(n)}\subseteq I_{n,K_n}.                   \tag{C}\label{eq:catcap}
\end{equation}
After fixing $n$ this is a bounded Boolean test, even though its validity
for every $n$ is a global arithmetic hypothesis. No monotonicity of $s$
is assumed.

\paragraph{Theorem.}
Condition \eqref{eq:catcap} at every $n$ is equivalent to
$I_{n,t}=B_n$ for all $n$ and $t\ge s(n)$.
From the former, a PR algorithm extends every given $a\in B_n$ to a
coherent sequence on all levels $n,n+1,\ldots$. The original bonding
maps may fail to be surjective.

\paragraph{Proof.}
If the images stabilize by $s(n)$, take $t=K_n$. Conversely, fix
$a\in B_n$. Condition \eqref{eq:catcap} yields a successful search
for $z\in A_{K_n}$ with $P_{n,K_n}(z)=a$. Put
$c=P_{n+1,K_n}(z)$. Projecting $z$ first to $s(n+1)$ shows that
$c\in B_{n+1}$, and $p_n(c)=a$ by the composite law.
Thus bounded search constructs a PR successor $L(n,a)$ with a
default value on invalid input. Define $w(0)=a$ and
$w(j+1)=L(n+j,w(j))$. At every $j$,
\[
     w(j)\in B_{n+j}, \qquad P_{n,n+j}(w(j))=a .
\]
Induct on the conjunction, compiling the first membership into its bounded
Boolean image test. The successor step is the lift just proved and one
composition of projections. In particular $p_{n+j}(w(j+1))=w(j)$.
For $t\ge s(n)$ choose $j=t-n$. Then $w(j)\in A_t$ supplies the missing
witness that $a\in I_{n,t}$. The reverse inclusion follows from decreasing
images. This proves the equivalence and exhibits the sequence as a
specified PR function of $n,a,j$. To see that $p_n(B_{n+1})\subseteq B_n$, apply the now proved stabilization at level $n+1$ and time $K_n$,
then project from $K_n$ to $s(n)$. The restricted tower is surjective.

\paragraph{Application to exactness.}
Let $q_n:G_n\to H_n$ be surjective homomorphisms of PR-presented
finite abelian groups, compatible with the bonds, and assume
\eqref{eq:catcap} with rate $s$ for the tower of kernels.
Every specified PR coherent $h_n\in H_n$ then has a PR coherent lift
$g_n\in G_n$. Indeed let $E_n=q_n^{-1}(h_n)$, a nonempty decidable fiber
with the inherited bonds. For $a\in I^E_{n,s(n)}$ take a preimage
$e\in E_{s(n)}$, and choose any $z\in E_{K_n}$ by bounded search.
The difference $e-P_{s(n),K_n}(z)$ lies in $\ker q_{s(n)}$.
Its image $a-P_{n,K_n}(z)$ at $n$ therefore lifts, by the kernel
comparison, to $d\in\ker q_{K_n}$. Now $z+d\in E_{K_n}$ projects to
$a$. This proves the comparison for $E$ with the very same rate.
Project any member of $E_{s(0)}$ to level zero for a seed, and apply the
walk just constructed. This gives an arithmetic, parameter-wise form of
exactness of the inverse limit at the right-hand term. All choices in
the argument are bounded scans.

For example, let every $A_n=\{0,1,2\}$ and every bond fix $0,1$ and
send $2$ to $1$. Then $s(n)=n+1$ works and $B_n=\{0,1\}$.
Even two elements per level permit towers for which no rate can be computed uniformly from a tower index.
For program $e$, let $A_0=\{0,1\}$. For $n>0$ use the same two elements
until bounded simulation by time $n$ reports a halt, then use only $0$.
The maps are inclusions. A uniform computable rate would decide halting since $1\in I_{0,s(0)}$ exactly when $e$ never halts, by stability and the
definition of the tower. This is an external impossibility result.
Every tower in this example already has the constant zero branch. The
information a rate recovers is extendibility of \emph{all} persistent
states, which is stronger than a single branch.

\paragraph{Lean endpoints.}
\path{lean/SheafCase.lean} proves the principal-cover gluing and both inverse equations in ordinary Lean, using the specialization preorder $x\preceq y$ iff $y\in N_x$. Its \texttt{SheafGluing} structure records the glue, its restrictions and uniqueness. The preceding proof constructs these operations from a finite sheaf table. \path{lean/CategoryCase.lean} proves the semantic certificate calculation
for the image lift and coherent iteration. In the object calculus implement
$P$ by the downward register loop and prove the split equation by the
resumed-run invariant above. Compile image membership and
\eqref{eq:catcap} with bounded reflection. The lift proof then supplies
\texttt{Tower.Lifting} for the new membership test $B_n$, so the existing
\texttt{CoherentTower.lean} route supplies \texttt{tower\_closed}.
For sheafification, implement masks and mixed-radix filtered tuples, prove
that decoding after projection deletes exactly the complementary coordinates,
and prove minimal-neighborhood membership by intersection-prefix induction.
The proof of the reflector then uses bounded operations for gluing tuple
entries, checking matching on intersections, and selecting the unique glue
in a supplied sheaf. Prove the two inverse equations pointwise and close
them by bounded reflection and table extensionality.
\texttt{FiniteLimits.lean} and \texttt{FiniteLimitTest.lean} already handle
fixed diagram shapes. Varying the shape requires a numerical scan over
its arrow table. For zero radices, separate an empty coordinate range
before invoking positive-radix digit lemmas.
